% !TeX root = ../colloquium-2.tex

% Определение локального экстремума и формулировка теоремы Ферма.

\phantomsection\label{sec_Definition11}
\begin{definition}[Определение локального экстремума и формулировка теоремы Ферма.]
    \subsection*{Определение локального экстремума.}
    \begin{note}[purple]
        Пусть $\delta > 0$, $U_{\delta}(a)$ - $\delta$-окрестность точки $a$. \\
        Если функция $f$ определена на множестве $E$, точка $a \in E$ и при всех 
        $x \in U_{\delta}(a) \cap E$, выполнено неравенство $f(x) \leq f(a)$, 
        то $a$ называется точкой локального минимума функции $f$. \\
        Если неравенство строгое, то $a$ называется точкой строгого локального максимума
        или строго локального минимума соответственно. Точки \textbf{максимума} и \textbf{минимума} называются
        точками \textbf{локального экстремума} (если неравенства нестрогие), и \textbf{точками строгого 
        локального экстремума}.
    \end{note}

    \subsection*{Теорема Ферма.}
    \begin{note}[purple]
        Пусть функция $f$ определена на интервале $(a, b)$, дифференцируема в точке $c \in (a, b)$
        и имеет в точке $c$ локальный экстремум. \\
        Тогда $f'(c) = 0$. 
    \end{note}
\end{definition}